\chapter{Security Foundations CIA Triad}
\label{ch:2.1}
\chaptertopic{Topic 2.1}

\chapterquote{Distrust and caution are the parents of security.}{Benjamin Franklin, \emph{Poor Richard's Almanack} (1758)}

\section{Motivation}

Unit 1 walked through a series of specific attacks -- social engineering, credential guessing,
wireless interception, AI-augmented deception. This chapter steps back and asks a more basic
question: what, exactly, is a security goal trying to protect, and how do you decide which risks
deserve your limited attention? Every security decision, no matter how specific, traces back to
three questions: is this being kept from people who shouldn't see it, has it been changed
without authorization, and can the people who need it actually reach it? Answering those
questions well also takes a shared vocabulary for deciding whether a given risk is worth acting on
at all. This chapter builds both:
the CIA triad as the vocabulary for \emph{what} is at stake, threat-actor categories as the
vocabulary for \emph{who} is attacking, and risk scoring as the vocabulary for \emph{how much} to
care.

\begin{apalignment}{Unit 2: Securing Spaces}{Topic 2.1 -- Cyber Foundations}
\begin{apobjectives}
\apobjective{2.1.A}{Analyze Risk}{Identify which of the three CIA-triad pillars -- confidentiality, integrity, or availability -- a given incident violates, including incidents that violate more than one.}
\apobjective{2.1.B}{Analyze Risk}{Distinguish the categories of adversary an organization may face, by the resources they command, the outcome they are pursuing, and the level of exposure they are willing to accept.}
\apobjective{2.1.C}{Analyze Risk}{Set out the stages an intrusion moves through, from the survey of a target to the actions the adversary came to perform, and explain why interrupting any one stage frustrates the whole sequence.}
\apobjective{2.1.D}{Analyze Risk}{Describe how an organization conducts a risk assessment: inventory the assets, identify the threats and weaknesses that bear on them, and estimate how likely a loss is and how large it would be.}
\apobjective{2.1.E}{Mitigate Risk}{Identify the responses available once a risk is understood -- reducing it, moving its financial consequence elsewhere, abandoning the activity that creates it, or accepting it as a recorded decision.}
\apobjective{2.1.F}{Mitigate Risk}{Sort security controls both by what they do in the timeline of an incident and by the means through which they are implemented.}
\apobjective{2.1.G}{Mitigate Risk}{Account for why layered, independent defenses are required rather than a single strong one, given that every individual control has a known failure mode.}
\end{apobjectives}
\end{apalignment}

\section{The CIA Triad}

\begin{definition}{CIA triad}
The three foundational security properties an organization protects: \textbf{Confidentiality},
\textbf{Integrity}, and \textbf{Availability}, formalized in NIST's core security publications as
the entry point for all risk analysis.
\end{definition}

\textbf{Confidentiality} is the property that information is accessible only to those authorized
to have it. Its threat is unauthorized disclosure -- an intercepted unencrypted login, a
misconfigured cloud bucket exposing customer records. Its controls are correspondingly
about restricting reach: \textbf{encryption} (transforming data so it is unreadable without a
key), \textbf{access control} (authenticating identity and limiting what each account can touch),
and classification policies that label data by sensitivity so handling rules apply automatically.
When confidentiality fails, data has been exposed, sometimes for months before anyone notices.

\textbf{Integrity} is the property that information and systems have not been altered without
authorization. Its threat is quiet modification -- a financial record changed, a software update
replaced with malware. Its signature control is the \textbf{cryptographic hash}: a
fixed-length fingerprint of data such that changing even one input bit changes the output
completely. A published software hash lets you verify what you downloaded matches what the
publisher released; a mismatch means tampering or corruption. Audit logs and version control
serve the same goal by creating an immutable record of who changed what, so tampering can be
detected and reversed. Integrity failures are insidious precisely because the system keeps
running -- it just quietly produces wrong answers.

\textbf{Availability} is the property that systems and data are reachable by authorized users when
needed. Its threats are denial in one form or another -- ransomware encrypting files until
payment, a distributed denial-of-service (\textbf{DDoS}) attack flooding a server until it
crashes, plain hardware failure. Its controls are about \textbf{redundancy}: backup
servers, mirrored data centers, and DDoS-mitigation traffic filtering. Disaster-recovery plans add
two more numbers to the mix: a Recovery Time Objective (how fast systems must come back) and a
Recovery Point Objective (how much data loss is tolerable). An e-commerce site down for a day
loses revenue and trust just as surely as a stolen customer list would.

\begin{example}[The triad in conflict]
The three pillars do not simply cooperate -- they routinely trade off against each other. Strong
encryption protects confidentiality, but adds computational overhead and key-management
complexity. If a key is lost or the decryption system fails, the data becomes unreachable -- an
availability failure created while solving a confidentiality problem. Strict access control
protects confidentiality just as directly, but can mean an emergency-room doctor cannot quickly
pull up a patient's drug allergies. That is a life-threatening availability failure, caused by a
confidentiality control working exactly as designed. Reasoning through exactly this kind of
tradeoff, for a specific context, is what the AP framework calls the Analyze Risk skill.
\end{example}

A useful habit when reading about any real incident is to ask which pillar or pillars were
violated. The 2017 WannaCry ransomware worm primarily hit \textbf{availability} -- it encrypted
files and demanded payment, locking hospitals out of patient records, without needing to read any
of the data it encrypted. A phishing campaign harvesting passwords targets
\textbf{confidentiality}. An attacker quietly editing a database entry to redirect payments hits
\textbf{integrity}. Many real incidents chain all three: breach confidentiality to steal
credentials, then use them to compromise integrity by modifying records. The final step is often
an availability attack -- crashing logging systems to cover tracks -- used purely to protect the
earlier two.

\begin{lstlisting}[language=Python, caption={Classifying which CIA pillars an incident violates}]
CONFIDENTIALITY_WORDS = {"leaked", "exposed", "stolen", "intercepted", "disclosed"}
INTEGRITY_WORDS = {"modified", "altered", "tampered", "corrupted", "falsified"}
AVAILABILITY_WORDS = {"encrypted", "ransom", "offline", "ddos", "crashed", "unavailable"}

def classify_incident(description):
    """Returns a sorted list of CIA pillars ("confidentiality", "integrity",
    "availability") whose keyword set appears in `description`."""
    words = set(description.lower().split())
    pillars = []
    if words & CONFIDENTIALITY_WORDS:
        pillars.append("confidentiality")
    if words & INTEGRITY_WORDS:
        pillars.append("integrity")
    if words & AVAILABILITY_WORDS:
        pillars.append("availability")
    return sorted(pillars)
\end{lstlisting}

\begin{misconception}
\textbf{``Security is mainly about keeping things secret.''} Confidentiality is only one of three
pillars; attacks on integrity or availability routinely cause equal or greater harm without
exposing a single record.
\end{misconception}

\section{Threat Actors and Motivation}

Knowing \emph{what} is at stake is only half of risk analysis. The other half is knowing
\emph{who} is likely attacking it, because a defense tuned for one actor type can be nearly
useless against another.

\begin{example}[Three tickets, three attackers]
A university security team arrives Monday to three incidents at once. The registrar's database
is encrypted with a ransom note attached. An anonymous forum post contains thousands of leaked
student records. And the athletics department website has been replaced with a political banner.
Three attacks in one morning, and three entirely different attackers behind them -- treating all
three the same would waste the response to each.
\end{example}

Security practitioners sort attackers into five recurring categories. \textbf{Nation-state actors}
are government-sponsored or -directed, pursuing geopolitical goals -- espionage, sabotage,
intelligence gathering -- with large budgets, skilled operators, and patience measured in years.
They are rarely prosecuted and rarely stop absent a diplomatic shift. \textbf{Cybercriminals} are
motivated purely by money, most visibly through ransomware -- software that encrypts a victim's
files and demands payment for the key -- often rented as ransomware-as-a-service from specialized
developers. Hospitals, universities, and municipalities are frequent targets precisely because
downtime is catastrophic for them, and because they are statistically more likely to pay.
\textbf{Hacktivists} combine hacking with activism, wanting embarrassment or amplification rather
than money -- a website defacement costs the target almost nothing technically but generates
outsized press. Because publicity \emph{is} the goal, hacktivist attacks tend to be loud and
openly attributed: the opposite of nation-state stealth.

\begin{definition}{Insider threat}
A threat originating from someone who already holds legitimate access to an organization's
systems. It spans three subtypes: malicious insiders (intentional theft or sabotage), disgruntled
employees (impulsive action after termination or a bad review), and negligent insiders (accidental
harm such as a misdirected email, a reused weak password, or an unvetted USB drive).
\end{definition}

Insider threats are uniquely hard to stop because the actor already has the credentials and
physical access every outside attacker has to fight for. Technical perimeter controls do nothing
against someone already inside the perimeter. Least-privilege access, user-activity monitoring,
and a rigorous off-boarding checklist are the controls that actually apply. \textbf{Script
kiddies} sit at the opposite end of sophistication -- low-skill attackers running public tools
they do not understand, motivated by curiosity or peer notoriety, scanning broadly for whatever
responds rather than targeting anyone specifically. They are the most common attacker by sheer
volume, and the easiest to stop: patching and closing unnecessary services eliminates most
of what they can do.

\begin{definition}{Advanced Persistent Threat (APT)}
A campaign \emph{style}, not a distinct actor type: a long-running, stealthy, high-capability
intrusion in which an attacker gains an initial foothold and then moves quietly through a network
for months or years, exfiltrating data while actively avoiding detection. It is almost always
nation-state-sponsored or nation-state-adjacent.
\end{definition}

Motivation predicts the highest-leverage defense, and mismatching the two wastes budget. Network
segmentation and anomaly detection answer a nation-state or APT threat. Offline backups, email
filtering, and MFA answer a cybercriminal's ransomware. A web application firewall and CMS
credential hygiene answer a hacktivist's defacement. Least privilege and access logging answer an
insider. Routine patching answers a script kiddie. None of these substitute for each other --
an organization facing all five actor types, as most do, needs all five defenses simultaneously.

\begin{misconception}
\textbf{``APT always means nation-state.''} APT describes the campaign's style -- advanced and
persistent -- not who is funding it. Well-resourced criminal groups can and do run APT-style
campaigns without any government sponsorship.
\end{misconception}

\begin{definition}{Cyber Kill Chain}
A model published by Hutchins, Cloppert, and Amin at Lockheed Martin in 2011\srccite{2.1}{1} that
breaks an intrusion into seven sequential, dependent stages: reconnaissance, weaponization,
delivery, exploitation, installation, command and control, and actions on objectives.
\end{definition}

The kill chain's usefulness rests on one structural claim: because each stage depends on the one
before it, interrupting the adversary at \emph{any single stage} denies them everything that would
have followed. A defender does not need to catch every step -- only one. The same dependency
structure holds for a purely physical intrusion, which is exactly why this unit keeps returning to
layered defense. An intruder targeting a building moves through the identical seven stages.
Reconnaissance means watching shift changes and photographing a staff badge. Weaponization means
cloning the badge and assembling a contractor pretext. Delivery means arriving at a loading dock
inside a normal delivery window. Exploitation means tailgating through a badged door. Installation
means concealing a rogue network device in a wiring closet. Command and control means that device
dialing out over a cellular link to bypass the firewall. And actions on objectives means copying
data from an unattended console. A staffed dock with booked deliveries, a mantrap admitting one
person at a time, and a locked, contact-alarmed wiring closet are candidate interruption points at
three different stages. Any one of the three, on its own, ends the intrusion.

\begin{figure}[h]
\centering
\begin{tikzpicture}[
  stage/.style={draw=kickstartermaroon, thick, rounded corners, minimum width=1.55cm, minimum height=1.3cm, align=center, font=\scriptsize, fill=kickstartercream},
  node distance=0.15cm
]
\node[stage] (s1) at (0,0) {Recon-\\naissance};
\node[stage] (s2) at (1.7,0) {Weapon-\\ization};
\node[stage] (s3) at (3.4,0) {Delivery};
\node[stage] (s4) at (5.1,0) {Exploit-\\ation};
\node[stage] (s5) at (6.8,0) {Install-\\ation};
\node[stage] (s6) at (8.5,0) {Command\\\& Control};
\node[stage] (s7) at (10.2,0) {Actions on\\Objectives};
\foreach \a/\b in {s1/s2, s2/s3, s3/s4, s4/s5, s5/s6, s6/s7}
  \draw[-{Latex[length=2mm]}, kickstartermaroon] (\a) -- (\b);
\node[font=\scriptsize\sffamily, gray, align=center] at (5.1,-1.2)
  {interrupting \textbf{any one} stage denies the attacker every stage that would have followed};
\end{tikzpicture}
\caption{The Cyber Kill Chain: seven sequential, dependent stages. Because each stage depends on
the one before it, a defender needs to reliably interrupt only one to end the intrusion.}
\label{fig:2.1-kill-chain}
\end{figure}

MITRE ATT\&CK is a complementary catalog rather than a competing model. Instead of a linear
sequence, it organizes observed adversary behavior into \emph{tactics} -- the objective at a given
moment, such as Initial Access or Persistence. Beneath each tactic sit the specific
\emph{techniques} used to pursue it, along with documented detections and mitigations. An
adversary in ATT\&CK may revisit a tactic, skip one entirely, or pursue two at once -- closer to
how real intrusions unfold than the kill chain's strict ordering. The kill chain is the better tool
for arguing that layered controls work at all. ATT\&CK is the better working reference once you
need to know what one specific technique looks like in a log or on a badge reader.

\section{The Language of Risk}

Once you can name what is threatened and by whom, the remaining question is how much any given
risk deserves your attention -- because no organization can defend against everything equally.

\begin{definition}{Risk}
The probability that a threat will exploit a vulnerability and cause harm to an asset. A
\textbf{threat} is the potential negative event itself (a hard drive failing, a database leak). A
\textbf{vulnerability} is the weakness that makes it possible (an outdated OS, a server room in a
flood-prone basement). Their combination realized is a security incident.
\end{definition}

NIST SP 800-30 formalizes the vocabulary used to quantify risk. \textbf{Asset value} ($V$) is the
monetary or strategic worth of what is being protected. \textbf{Likelihood} ($L$) is the
probability of the threat occurring, scored 1 (very low) to 5 (very high). \textbf{Impact} ($I$) is
the severity of the damage if it does occur, scored 1 (negligible) to 5 (catastrophic). These
combine into a single number:
\[
\text{Risk Score} = L \times I,
\]
ranging from 1 to 25, which lets an organization rank many risks against each other on one scale
rather than arguing about each in isolation.

A computed score then resolves into one of four canonical treatments. \textbf{Mitigate} reduces
likelihood or impact through a control -- a fire-suppression system reduces impact, regular
vulnerability scanning reduces likelihood -- and is the right call whenever a cost-effective
control exists for a high-scoring risk. \textbf{Accept} takes no action, appropriate when a risk
score is low or mitigation costs exceed the asset's value. A company might accept that its
low-traffic marketing blog could go down for an hour, since the blog carries no transactional
value. \textbf{Transfer} shifts the financial burden to a third party -- cyber insurance, or a
cloud host's SLA that pays out during downtime -- fitting high-impact risks too expensive to
mitigate fully but insurable. \textbf{Avoid} eliminates the risk by changing the underlying
activity entirely. A retailer that stops storing customer card numbers and routes payments through
a third-party processor has not mitigated the risk of a card-database breach -- it has removed the
database.

\begin{lstlisting}[language=Python, caption={Recommending a risk treatment from likelihood, impact, and cost}]
def recommend_treatment(likelihood, impact, mitigation_cost, asset_value):
    """likelihood, impact: ints 1-5. mitigation_cost, asset_value: numbers.
       Returns one of "Mitigate", "Accept", "Transfer", "Avoid"."""
    risk_score = likelihood * impact
    if risk_score <= 4:
        return "Accept"
    if mitigation_cost > asset_value:
        return "Transfer" if impact >= 4 else "Avoid"
    return "Mitigate"
\end{lstlisting}

The comparison that ultimately decides between these four is always the same shape: weigh the
cost of a control against the asset's value and the risk score together. Spending \$20{,}000 to
protect a \$10{,}000 database is mathematically irrational, regardless of how the risk score reads
in isolation. Risk language exists to make that comparison explicit rather than left to
intuition.

\subsection{Sorting the Controls Themselves}

Mitigation is a category of treatment, not a single action -- and once mitigation is the chosen
path, a control still needs to be classified two more ways: when in an incident's timeline it
acts, and what kind of mechanism it is.

By \textbf{timeline}, a control is \textbf{preventive} if it stops an incident before it starts (a
locked door, input validation, MFA), \textbf{detective} if it identifies an incident already in
progress or already completed (an intrusion-detection alert, a smoke detector, an audit log
review), or \textbf{corrective} if it limits damage or restores normal operation after the fact (a
backup restore, an incident-response runbook, a patch rolled out after a vulnerability is
exploited). By \textbf{means}, a control is \textbf{technical} if it is enforced by a system
(encryption, a firewall rule, a badge reader), \textbf{administrative} if it is enforced by a
policy or a process (a background-check requirement, a mandatory security-awareness training,
separation-of-duties in an approval workflow), or \textbf{physical} if it acts on the physical
world directly (a fence, a security guard, a locked server-room door).

The two axes are independent, so the same control can occupy exactly one cell of a $3\times3$
grid. A badge reader is technical and preventive. A security guard reviewing camera footage after
an incident is physical and detective. A mandatory two-person sign-off before a wire transfer is
administrative and preventive. Naming both axes for a proposed control is what keeps a security
plan from silently duplicating coverage in one cell -- three preventive technical controls -- while
leaving another cell -- any detective administrative control -- completely empty.

\begin{misconception}
\textbf{``We must eliminate all risk.''} This is not achievable, and attempting it is itself a
cost. Security has a price, and every organization operates with some accepted level of residual
risk -- its risk appetite -- in order to function at all.
\end{misconception}

\section*{Review Questions}
\begin{reviewquestions}
\item An attacker steals a database of customer emails and, separately, silently edits shipping
addresses in an order-processing system to redirect packages. Identify which CIA pillar each
action violates and explain why they require different remediation steps.
\item Explain why a strict access-control policy, which strengthens confidentiality, can create an
availability failure in an emergency-medicine context, and describe how a security architect might
balance the two.
\item A university discovers a complete, cleanly organized set of student records posted publicly.
Explain why the completeness and cleanliness of the leak points toward an insider threat rather
than an external attacker, and name two controls that specifically address insider risk.
\item A risk is scored at likelihood 2, impact 5. Using the risk-score formula and the four
treatment categories from this chapter, explain which treatment is most appropriate and what
additional information (beyond the score itself) would change that recommendation.
\item Explain the difference between a threat, a vulnerability, and a risk using a concrete
example, and explain why conflating "risk" with "threat" leads to poor prioritization decisions.
\item A physical intrusion is stopped at the loading dock because deliveries must be pre-booked
and verified, before the intruder ever reaches a badged door. Using the Cyber Kill Chain's
structural claim about dependent stages, explain why this single control is sufficient even
though the building still has no mantrap and an unlocked wiring closet.
\item MITRE ATT\&CK is described as a better working reference than the Cyber Kill Chain for
identifying what a specific logged event represents, even though the Kill Chain is the better tool
for justifying layered defense in the first place. Explain why.
\item Classify each control by both timeline (preventive, detective, or corrective) and means
(technical, administrative, or physical): (a) a fire-suppression system, (b) a mandatory
background check before hiring, (c) reviewing badge-access logs the morning after a break-in.
Then explain what it would mean for a security plan if every proposed control landed in the same
cell of the grid.
\end{reviewquestions}

\begin{answerkey}
\item The theft of the email database is a confidentiality failure: information reached parties
never authorized to see it. The remediation is disclosure notification, credential resets for
affected customers, and an audit of how the data was exposed. The silent edit to shipping addresses
is an integrity failure: authorized data was altered without authorization. The remediation is
restoring from a trusted backup, auditing who had write access, and adding change-detection
controls such as hashing or versioning on order records. The two failures require different
remediation. Confidentiality damage cannot be undone once data is copied -- the response is
entirely forward-looking. Integrity damage, by contrast, can often be reversed once detected, but
only if a trustworthy prior state exists to restore from.
\answerlabel{2.1.A}
\item Strict access control strengthens confidentiality by narrowing who can reach a patient
record. But in an emergency-medicine context, that same narrowing can block a treating clinician
from information needed within seconds -- an availability failure caused by the confidentiality
control working exactly as designed. A security architect balances the two with
context-aware controls rather than a single static policy. Break-glass access grants an
on-call clinician immediate access while logging the override for mandatory after-the-fact review.
Role-based access stays scoped tightly in ordinary circumstances but expands under a documented
emergency procedure. Monitoring treats every break-glass event as an audit item rather
than a routine action.
\answerlabel{2.1.G}
\item A large, external, opportunistic leak is typically noisy, partial, and littered with
extraction artifacts. That is because an external attacker is working through whatever access they
managed to obtain, exfiltrating in the shape that access allowed. A complete, cleanly organized
dataset instead looks like a direct, authorized export -- the kind only someone with legitimate,
well-formed access to the underlying system could produce with minimal effort, which points toward
an insider. Two controls specifically address insider risk. Least-privilege access ensures no
single account, insider or not, can export the entire dataset in one legitimate query.
User-activity monitoring and data-loss-prevention alerting on unusually large exports flags
the action itself, rather than relying on detecting external intrusion.
\answerlabel{2.1.B}
\item Risk score $= L \times I = 2 \times 5 = 10$, a moderate-to-high score driven entirely by
impact rather than frequency. Given a rare-but-catastrophic profile, mitigation is often
appropriate if an effective, reasonably priced control exists -- an impact-5 event is by definition
one an organization cannot absorb casually. But transfer is equally defensible if impact is high
enough that even a well-mitigated residual loss is worth insuring, and mitigation is not
cost-effective. The score alone cannot decide between mitigate, transfer, or avoid: what is missing
is the mitigation cost against asset value, the availability of insurance for this specific
risk, and whether the underlying activity can be eliminated (avoidance) without sacrificing a
function the organization actually needs.
\answerlabel{2.1.E}
\item A threat is the potential negative event itself -- a fire, a stolen credential, a flood. A
vulnerability is the specific weakness that would let that event cause harm -- an unsuppressed
server room, a password with no expiration policy, a basement data center. Risk is the probability
that the threat exploits that vulnerability, combined with the resulting harm; it is only
meaningful as a combination of the two. Conflating "risk" with "threat" leads an organization to
react to whatever threat is most recently in the news (chasing the threat) rather than to where a
real vulnerability and a real asset intersect. That routinely misallocates a security budget toward
exotic, unlikely threats, while an ordinary, exploitable weakness with no threat attached at the
moment goes unpatched until a threat does show up.
\answerlabel{2.1.D}
\item The kill chain's stages are dependent: weaponization requires reconnaissance to have
succeeded, delivery requires a weapon, exploitation requires successful delivery, and so on
through every later stage. Stopping the intruder at delivery -- the loading dock -- denies them
every stage that would have followed, including exploitation, installation, and actions on
objectives. That holds regardless of whether those later stages would themselves have been easy or
hard to execute. A missing mantrap or an unlocked wiring closet are real weaknesses, but they are
weaknesses in stages the intruder never reaches, so they impose no cost on this particular
attempt.
\answerlabel{2.1.C}
\item The kill chain enforces a strict linear order that real intrusions do not always follow --
an adversary may skip a stage, revisit one, or pursue two objectives at once. So mapping one
specific logged event (an unusual outbound connection, a suspicious registry change) onto a rigid
seven-stage sequence can be awkward or ambiguous. ATT\&CK instead catalogs tactics and, beneath
each, the specific documented techniques adversaries use to pursue them, along with known
detections and mitigations. That makes it the more direct reference when the question is "what
does this one observed behavior mean," rather than "why does interrupting any single stage work."
\answerlabel{2.1.C}
\item (a) A fire-suppression system is technical and corrective -- it is enforced by a system, and
it acts after ignition has already occurred, to limit damage rather than to prevent it. (b) A
mandatory background check is administrative and preventive -- it is a policy requirement, applied
before the hire, meant to stop a bad outcome from ever starting. (c) Reviewing badge-access logs
the morning after a break-in is administrative and detective -- it is a process, not a system
enforcing anything automatically, and it identifies what already happened rather than preventing
or repairing it. If every proposed control landed in the same cell, the security plan would be
lopsided in a specific, predictable way: strong on paper wherever that one cell's strengths lie,
and silently unguarded everywhere else -- for example, an organization with only preventive
technical controls has nothing in place to detect a control that has already failed, and nothing
to recover with once one does.
\answerlabel{2.1.F}
\end{answerkey}

\section*{Sources}
\begin{enumerate}[label=\lbrack\arabic*\rbrack, leftmargin=2.2em]
\item \srcitem{2.1}{1}Eric M. Hutchins, Michael J. Cloppert, and Rohan M. Amin, ``Intelligence-Driven
Computer Network Defense Informed by Analysis of Adversary Campaigns and Intrusion Kill Chains,''
Lockheed Martin Corporation, \emph{Leading Issues in Information Warfare \& Security Research},
2011.
\end{enumerate}
